\documentclass[11pt]{article}
\usepackage{fontspec}
\usepackage{amsmath,amssymb,amsthm}
\usepackage{fancyvrb}
\usepackage{booktabs}
\usepackage{stmaryrd}
\usepackage{tikz-cd}
\usepackage[margin=2.7cm]{geometry}
\usepackage{xcolor}
\usepackage[colorlinks,linkcolor=black,citecolor=black,urlcolor=blue!60!black]{hyperref}

% fonts are loaded by file name (as required by arXiv's XeLaTeX setup)
\setmainfont{lmroman10-regular.otf}[BoldFont=lmroman10-bold.otf,ItalicFont=lmroman10-italic.otf,
  BoldItalicFont=lmroman10-bolditalic.otf,SlantedFont=lmromanslant10-regular.otf]
\setsansfont{lmsans10-regular.otf}[BoldFont=lmsans10-bold.otf,ItalicFont=lmsans10-oblique.otf,
  BoldItalicFont=lmsans10-boldoblique.otf]
\setmonofont{JuliaMono-Regular.ttf}[BoldFont=JuliaMono-Bold.ttf,ItalicFont=JuliaMono-RegularItalic.ttf,
  BoldItalicFont=JuliaMono-BoldItalic.ttf,Scale=0.82]
\DefineVerbatimEnvironment{agda}{Verbatim}{xleftmargin=1.5em,fontsize=\small}
\newcommand{\ag}[1]{\texttt{\detokenize{#1}}}
\newcommand{\repourl}{https://github.com/txa/sst-new}
\newcommand{\agfile}[1]{\href{\repourl/blob/main/#1}{\ag{#1}}}
\newcommand{\agdir}[1]{\href{\repourl/tree/main/#1}{\ag{#1/}}}

% path composition, as in Agda
\newcommand{\pc}{\mathbin{\text{\ttfamily ∙}}}

\newtheorem*{remark}{Remark}

\title{Now available in grey\\[0.5ex]
\large Semi-simplicial groupoids from weakenings and thinnings}
\author{Thorsten Altenkirch\\[0.5ex]
\small University of Nottingham, United Kingdom\\
\small\href{mailto:Thorsten.Altenkirch@nottingham.ac.uk}{\ag{Thorsten.Altenkirch@nottingham.ac.uk}}}
\date{October 2026}

\begin{document}
\maketitle

\begin{abstract}
In a companion note~\cite{obstruction} we defined a small type theory of
weakenings and thinnings, in which semi-simplicial types up to level $n$ can
be constructed uniformly in $n$ as contexts $\mathrm{SST}\,n$. The syntax is
defined by induction-recursion without K, and it is a set. Its naive
interpretation in the universe, however, needs uniqueness of identity proofs
(UIP) for the semantic types. Here we lift this limitation by one level. We
interpret the syntax in families of \emph{groupoids}, without UIP and without
K. Two changes make this work: a \emph{Forded} semantics of telescope
extension, in which a kept entry never transports its value but only changes a
stored path, and a semantic action of the codes for equality of thinnings,
defined by recursion. What remains are 2-dimensional coherences of the
semantic composition and weakening, which we prove by induction; their
3-dimensional components are trivial because the semantic types are
groupoids. As a result, $\llbracket\mathrm{SST}\,n\rrbracket$ is a type of
semi-simplicial groupoids up to level $n$, defined uniformly in $n$. We
also report on the next level, 2-types, where the 3-dimensional coherences have
to be proved. The development is formalised in Cubical Agda, in the repository
\href{\repourl}{\ag{github.com/txa/sst-new}}.
\end{abstract}

\begin{quote}\small
``Any customer can have a car painted any colour that he wants so long as it
is black.''\par\hfill --- Henry Ford, \emph{My Life and Work}~\cite{ford}
\end{quote}
\noindent
From 1914 Ford's Model T was sold only in black; other colours returned in
1926. The companion note~\cite{obstruction}, whose title borrows Ford's
remark, found that the naive interpretation of our syntax works in any
colour, as long as it is UIP, that is, in sets. This note offers the next
shade: groupoids.

\section{Introduction}

Semi-simplicial types are the simplest structure in homotopy type theory that
needs an infinite tower of coherences, and it is a well-known open problem
whether the type of semi-simplicial types up to level $n$ can be defined
uniformly in a variable $n:\mathbb N$ in book HoTT~\cite{ufias-sst,shulman-eat}.
In two-level type theory this is possible~\cite{ack-strict,acks-2ltt}, and
Kolomatskaia and Shulman define them in displayed type theory~\cite{dtt}.

The companion note~\cite{obstruction} explains an attempt that proceeds in
two steps. First define a small type theory (\ag{TT}) whose contexts can
represent the data of semi-simplicial types, and construct the contexts
$\mathrm{SST}\,n$ uniformly in that syntax. Then interpret the syntax in the
universe. The first step goes through without K. The second one does not:
the naive interpretation transports along the semantic images of equations
that hold only propositionally, and making these transports agree needs UIP.
The note also explains that the problem cannot simply be moved into the
syntax: every presentation we tried either has critical pairs in the syntax or
a naive model that needs UIP.

This note shows that the limitation can be lifted, at least by one level.
We interpret \ag{TT} in families of groupoids, without UIP and without K
(file \agfile{experiments/sem1/Sem1.agda}). The coherences that the naive
interpretation takes from UIP are now proved by induction on the syntax. They
are 2-dimensional, and the only use of the truncation is that their
3-dimensional components are trivial in groupoids. As a consequence
\[
\llbracket\mathrm{SST}\,n\rrbracket : \mathcal U_1
\]
is a type of semi-simplicial groupoids up to level $n$, defined uniformly in
$n$. We do not claim that this is the first such definition; truncated
versions of the problem need only finitely many coherences, and are expected
to be solvable. The point is that the same syntax, and the same uniform
construction of $\mathrm{SST}\,n$, work, and that the semantics has a clear
structure that suggests how to go further.

Sections~\ref{sec:sst}--\ref{sec:naive} recall the syntax, the construction of
$\mathrm{SST}\,n$ and the naive interpretation from~\cite{obstruction}.
Section~\ref{sec:grpd} is the main part: the groupoid semantics.
Section~\ref{sec:sstgrpd} discusses the resulting semi-simplicial groupoids,
and Section~\ref{sec:higher} the next level.

\section{Semi-simplicial types up to triangles}\label{sec:sst}

A semi-simplicial type up to level 2 consists of
\begin{align*}
X_0 &: \mathcal U\\
X_1 &: X_0 \to X_0 \to \mathcal U\\
X_2 &: (x_0\,x_1\,x_2 : X_0) \to X_1\,x_0\,x_1 \to X_1\,x_0\,x_2 \to X_1\,x_1\,x_2 \to \mathcal U
\end{align*}
Points, edges between points, and triangles filling a boundary of three
points and three edges. Each $X_k$ is a family over a \emph{telescope}, and
each entry of the telescope is an earlier family applied to a
\emph{selection} of earlier entries: the edge $x_{02} : X_1\,x_0\,x_2$
selects $x_0$ and $x_2$. Selections compose, and at level 2 this composition
is used for the first time in a non-trivial way, since the boundary of a
triangle consists of edges and the boundary of an edge consists of points.

\section{The type theory of weakenings and thinnings}\label{sec:tt}

The syntax has five sorts (see~\cite[\S\S 3--5]{obstruction} for a more
leisurely account):
\begin{center}
\begin{tabular}{lll}
\toprule
\ag{Con : Set} & contexts & sequences of type families $X_0, X_1, \dots$\\
\ag{Ty : Con → Set} & types & telescopes (domains of families)\\
\ag{Var : (Γ : Con) → Ty Γ → Set} & variables & a family in $\Gamma$ over a telescope\\
\ag{Wk : Con → Con → Set} & weakenings & here only \ag{p : Wk (Γ ▷ A) Γ}\\
\ag{Thin : Ty Γ → Ty Γ → Set} & thinnings & projections between telescopes\\
\bottomrule
\end{tabular}
\end{center}
A thinning keeps some entries of a telescope and drops the others; read
backwards, it is an order-preserving injection of the entries of its target
into those of its source. The constructors are
\begin{agda}
•    : Con
_▷_  : (Γ : Con) → Ty Γ → Con          -- add a family over a telescope
ε    : Ty Γ                            -- empty telescope
ext  : (A : Ty Γ) → Var Γ C → Thin A C → Ty Γ
                                       -- A extended by an entry  x (δ a)
zero : Var (Γ ▷ A) (wkTy A p)
suc  : Var Γ A → Var (Γ ▷ B) (wkTy A p)
tε   : Thin ε ε
t-d  : Thin A B → Thin (ext A x γ) B   -- drop the last entry
t-k  : (σ : Thin A B) → CodeV x y → Code (δ ∘ σ) δ'
     → Thin (ext A x δ') (ext B y δ)   -- keep it
\end{agda}
Weakening and composition are operations,
\begin{agda}[frame=single,framesep=2mm]
wkTy   : Ty Γ → Wk Δ Γ → Ty Δ
wkVar  : Var Γ A → (w : Wk Δ Γ) → Var Δ (wkTy A w)
wkThin : Thin A B → (w : Wk Δ Γ) → Thin (wkTy A w) (wkTy B w)
_∘_    : Thin B C → Thin A B → Thin A C
\end{agda}
defined by recursion, simultaneously with the syntax, since they occur in the
indices of \ag{zero}, \ag{suc} and \ag{t-k}. So the syntax is an
inductive-inductive-recursive definition. Agda accepts it, but we know of no
general account of definitions of this form; we discuss this briefly in
Section~\ref{sec:discussion}.

A kept entry stores the equation that relates its index \ag{δ'} to the
composite \ag{δ ∘ σ} (it is \emph{Forded}). Instead of an identity type it
uses a code for equality of thinnings, defined by recursion:
\begin{agda}
Code : Thin A C → Thin A' C' → Set
Code tε          tε          = ⊤
Code (t-d σ)     (t-d τ)     = Code σ τ
Code (t-k σ _ _) (t-k τ _ _) = Code σ τ
Code _           _           = ⊥
\end{agda}
and similarly \ag{CodeV} for variables. Codes are propositions, by a plain
induction, and they decode to equations (\ag{Code-sound}). With them
composition is defined without transport, for instance
\begin{agda}
t-k δ ey q ∘ t-k σ ex e = t-k (δ ∘ σ) (CodeV-trans ex ey) (k-code q e)
\end{agda}
and the syntax is K-free (file \agfile{TT.agda}). Its sorts are sets, by
encode-decode~\cite[Theorem~7.2.2]{hott-book}.

\paragraph{SST in the syntax.}
Following~\cite[Section~3]{ack-strict}, we define by recursion on $k$,
simultaneously, the context \ag{ctx k} of levels $X_0,\dots,X_{k-1}$, the
$k$-skeleton \ag{SK k n : Ty (ctx k)} of an $(n-1)$-simplex, its face maps
\ag{SK→ k f : Thin (SK k n) (SK k m)} for order preserving embeddings $f$, and
their functoriality \ag{α k f g}, a code relating \ag{SK→ k g ∘ SK→ k f} and
\ag{SK→ k (f ⊙ g)}. This gives
\begin{agda}
SST : ℕ → Con
\end{agda}
uniformly in $n$ and without K (file \agfile{SSTk.agda}).

\section{The naive interpretation needs UIP}\label{sec:naive}

The intended interpretation of a context is a nested $\Sigma$-type of
families, of a telescope a type depending on the context, and of a thinning a
projection:\footnote{We write $\Sigma$-types as \ag{(x : A) × B}, by analogy
with \ag{(x : A) → B}, and let \ag{×} associate to the right. This is not
Agda syntax: in Agda the type is written \ag{Σ A (λ x → B)}, or
\ag{Σ[ x ∈ A ] B} with the standard library.}
\begin{agda}
⟦ Γ ▷ A ⟧C       = (γ : ⟦ Γ ⟧C) × (⟦ A ⟧T γ → Set)
⟦ ext A x δ ⟧T γ = (a : ⟦ A ⟧T γ) × ⟦ x ⟧V γ (⟦ δ ⟧Th a)
⟦ t-k σ _ e ⟧Th (a , v) = ⟦ σ ⟧Th a , subst ⟦ x ⟧V (keq e a) v
\end{agda}
where \ag{keq e a : ⟦ δ' ⟧Th a ≡ ⟦ δ ⟧Th (⟦ σ ⟧Th a)} is built from the
interpretation of the stored code and from
\begin{agda}
⟦∘⟧ : ⟦ δ ∘ σ ⟧Th a ≡ ⟦ δ ⟧Th (⟦ σ ⟧Th a)
\end{agda}
which holds only propositionally. In the case \ag{t-k ∘ t-k} of \ag{⟦∘⟧}
itself, a value is transported along two different paths, and we have to
show that the results agree. Since the semantic types are arbitrary, this is
UIP for them~\cite[\S 7]{obstruction}.

Note that this is a property of the naive \emph{semantics}, not of the syntax.
The syntax is a set, and that does not need UIP: it follows from the codes by
encode-decode. Nor does set-ness of the syntax restrict its interpretations,
since the syntax is not a quotient. What fails is the particular way in which
the naive interpretation interprets a kept entry, by transporting its value
along the image of a propositional equation. With arbitrary types the
coherence of these transports cannot be proved, and it is what the
interpretation in~\cite{obstruction} takes from a UIP assumption
(\ag{pair-irr}, file \agfile{Sem.agda}).

\section{The groupoid semantics}\label{sec:grpd}

We now restrict the semantic families to groupoids and show that then all the
coherences can be proved. The interpretation is in Cubical
Agda~\cite{cubical-agda}, file \agfile{experiments/sem1/Sem1.agda}. For equality
reasoning we use Agda's inductive identity type, for which J computes; it is
equivalent to the cubical path type (\ag{Cubical.Data.Equality}), and we use
the cubical library only to state and transfer h-levels.

\subsection{The interpretation}

\paragraph{Groupoids.}
A family in a context is now a family of groupoids:
\begin{agda}
record Grpd : Set₁ where
  field ⌞_⌟ : Set ; isGrpd : isGroupoid ⌞_⌟

⟦ • ⟧C     = ⊤
⟦ Γ ▷ A ⟧C = (γ : ⟦ Γ ⟧C) × (⟦ A ⟧T γ → Grpd)
\end{agda}
The interpretations of types (telescopes) are then groupoids as well
(\ag{isGrpdT}), since groupoids are closed under $\Sigma$ and path types of
groupoids are sets. This is all we need: the stored paths, and the
coherences between them, live in the interpretations of telescopes. The
interpretation of a context is not a groupoid, since it contains the universe
of groupoids, which by univalence is a 2-type; but nothing is needed about
its h-level.

\paragraph{Forded extension.}
The first change is to Ford the semantics of telescope extension, just as the
syntax Fords the index of a kept entry:
\begin{agda}
⟦ ext A x δ ⟧T γ = (a : ⟦ A ⟧T γ) × (d : ⟦ C ⟧T γ) × (⟦ δ ⟧Th a ≡ d) × ⟦ x ⟧V γ d
\end{agda}
An entry is a quadruple $(a, d, q, v)$: the earlier entries $a$, a point $d$
of the telescope $C$ of the family $x$, a path $q$ from the selection
$\llbracket\delta\rrbracket\,a$ to $d$, and a value $v$ of the family at $d$.
By contracting the singleton $(d,q)$ this is equivalent to the naive
$\Sigma\,(a : \llbracket A\rrbracket\,\gamma)\,.\,\llbracket x\rrbracket\,\gamma\,(\llbracket\delta\rrbracket\,a)$.
The advantage is that a kept entry no longer has to transport its value; it
only extends the stored path:
\begin{agda}
⟦ t-k σ ex e ⟧Th (a , d , q , v) = ⟦ σ ⟧Th a , d , kp δ σ δ' e a ∙ q , v
  where kp δ σ δ' e a = sym (⟦∘⟧ δ σ a) ∙ ≈ e a
\end{agda}
(We have suppressed the code \ag{ex : CodeV x y} relating the source and
target variables; in the formalisation it is decoded to a path in the set of
variables and matched with \ag{refl}.) A thinning never changes the point
$d$ or the value $v$ of an entry. Consequently, a path between the
interpretations of two thinnings has no component in the values: it is a path
in the earlier entries together with a 2-path between stored paths.

\paragraph{Codes act as paths.}
The second change concerns the stored codes. We need a path from
$\llbracket\delta\circ\sigma\rrbracket\,a$ to
$\llbracket\delta'\rrbracket\,a$ for every code
\ag{e : Code (δ ∘ σ) δ'}. Rather than applying the interpretation to the
decoded equation \ag{Code-sound e}, which does not compute, we give codes a
semantic action, by recursion:
\begin{agda}
≈ : Code σ τ → (a : ⟦ A ⟧T γ) → ⟦ σ ⟧Th a ≡ ⟦ τ ⟧Th a
\end{agda}
On \ag{tε} and \ag{t-d} it is \ag{refl} and the recursive call. On two keeps
it is a path in the Forded $\Sigma$-type, whose components are \ag{≈} on the
kept thinnings and a 2-path between stored paths, described below.

\paragraph{Weakening.}
The variables \ag{zero : Var (Γ ▷ B) (wkTy B p)} and
\ag{suc x : Var (Γ ▷ B) (wkTy A p)} live over weakened telescopes. A telescope
\ag{A} of \ag{Γ}, moved into \ag{Γ ▷ B}, has the same shape, but its entries
refer to the families \ag{suc x} instead of \ag{x}. Its interpretation
\ag{⟦ wkTy A p ⟧T (γ , X)} is therefore not definitionally
\ag{⟦ A ⟧T γ}: the two agree only if the interpretations of
\ag{wkThin δ p} and \ag{δ} agree, and that holds only up to a path. So we do
not identify them but relate them by a map, defined by recursion on \ag{A},
\begin{agda}
unwk : ⟦ wkTy A p ⟧T (γ , X) → ⟦ A ⟧T γ
unwk ε b                   = tt
unwk (ext A x δ) (a , d , q , v)
  = unwk A a , unwk C d , sym (unwk-nat δ a) ∙ cong (unwk C) q , v
\end{agda}
together with its naturality for thinnings,
\begin{agda}
unwk-nat : unwk C (⟦ wkThin δ p ⟧Th a) ≡ ⟦ δ ⟧Th (unwk A a)
\end{agda}
The variables are interpreted through it,
\begin{agda}
⟦ zero  ⟧V (γ , X) b = ⌞ X (unwk B b) ⌟
⟦ suc x ⟧V (γ , X) b = ⟦ x ⟧V γ (unwk A b)
\end{agda}
which is what makes the value \ag{v} in the definition of \ag{unwk} well
typed. Again the Forded semantics pays off: \ag{unwk} leaves the value alone
and only corrects the stored path, using \ag{unwk-nat}.

\paragraph{Paths of kept entries.}
All paths between kept entries have the same form. We state it for a general
Forded $\Sigma$-type. Let \ag{B} be a type of earlier entries, \ag{D} a type
of points with a selection \ag{f : B → D}, and \ag{V} a family over \ag{D};
for \ag{⟦ ext A x δ ⟧T γ} these are \ag{⟦ A ⟧T γ}, \ag{⟦ C ⟧T γ},
\ag{⟦ δ ⟧Th} and \ag{⟦ x ⟧V γ}. Given earlier entries \ag{b b' : B}, a
point \ag{d : D}, stored paths \ag{q : f b ≡ d} and \ag{q' : f b' ≡ d}, and a
value \ag{v : V d}, we can form the two entries
\begin{agda}
(b , d , q , v) , (b' , d , q' , v) : (b : B) × (d : D) × (f b ≡ d) × V d
\end{agda}
which share their point and their value. Given moreover a path
\ag{r : b ≡ b'} between the earlier entries and a 2-path
\ag{s : sym (cong f r) ∙ q ≡ q'}, saying that going back along the image of
\ag{r} and then along \ag{q} gives \ag{q'}, we have
\begin{agda}
pathK r s : (b , d , q , v) ≡ (b' , d , q' , v)
\end{agda}
defined by J on \ag{r}.

\ag{pathK} organises everything that follows. It splits a path between kept
entries into a \emph{base} part \ag{r}, which will be given by recursion or
induction on the thinnings, and a \emph{stored-path} part \ag{s}, which is
where the coherences come in. One level up, two composites of \ag{pathK}s are
equal if their base paths are equal and their stored-path parts agree; the
latter is an equation between 2-paths, and this is the only place where we
will use that the semantic types are groupoids
(Section~\ref{sec:where}).

The interpretation is defined by mutual recursion: the
interpretation of thinnings \ag{⟦_⟧Th}, the action of codes \ag{≈}, the
composition path \ag{⟦∘⟧}, and \ag{unwk} and \ag{unwk-nat}.
In each of \ag{≈}, \ag{⟦∘⟧} and \ag{unwk-nat}, the case of a kept entry is a
\ag{pathK}, whose base path is a recursive call and whose second component is
one of three 2-cells:
\begin{center}
\begin{tabular}{lll}
\toprule
cell & used in & statement\\
\midrule
\ag{cohN} & \ag{≈} on keeps & stored paths are natural in codes\\
\ag{cohA} & \ag{⟦∘⟧} on two keeps & stored paths of two keeps compose to that of the composite\\
\ag{cohW} & \ag{unwk-nat} on keeps & weakening commutes with stored paths\\
\bottomrule
\end{tabular}
\end{center}
Their types say that certain diagrams of paths commute. We omit
$\llbracket-\rrbracket$, writing $\delta\sigma a$ for
$\llbracket\delta\rrbracket(\llbracket\sigma\rrbracket\,a)$, write
$\langle c\rangle$ for the code path \ag{≈ c a}, and write
$\mathrm{kp}_e$ for the stored path \ag{kp δ σ δ' e a} of a keep with code
\ag{e : Code (δ ∘ σ) δ'}.
\begin{center}
\begin{tikzcd}[column sep=huge,row sep=huge]
\delta\sigma a \ar[r,"\delta\langle c\rangle"] \ar[dr,"\mathrm{kp}_e"'] & \delta\tau a \ar[d,"\mathrm{kp}_{e'}"] \\
& \delta' a
\ar[from=1-1,to=2-2,phantom,"\text{\ag{cohN}}" pos=0.55,shift left=2.5ex]
\end{tikzcd}
\qquad
\begin{tikzcd}[column sep=large]
\delta_C((\delta\circ\sigma)a) \ar[r,"{\delta_C\circ(\delta,\sigma)}"] \ar[dr,"\mathrm{kp}_{e'}"'] & \delta_C\delta\sigma a \ar[r,"\mathrm{kp}_{q}"] & \delta_B\sigma a \ar[dl,"\mathrm{kp}_e"] \\
& \delta_A a &
\ar[from=1-2,to=2-2,phantom,"\text{\ag{cohA}}" pos=0.45]
\end{tikzcd}

\medskip
\begin{tikzcd}[column sep=large]
U(\delta^+\sigma^+a) \ar[r,"{u(\delta)}"] \ar[d,"U\mathrm{kp}_{e_1}"'] & \delta\,U(\sigma^+a) \ar[r,"{\delta\,u(\sigma)}"] & \delta\sigma\bar a \ar[d,"\mathrm{kp}_e"] \\
U(\delta'^+a) \ar[rr,"{u(\delta')}"'] & & \delta'\bar a
\ar[from=1-1,to=2-3,phantom,"\text{\ag{cohW}}"]
\end{tikzcd}
\end{center}
\begin{itemize}
\item \ag{cohN}: for \ag{c : Code σ τ}, moving along the code path
  of the kept thinnings and then using the stored path of the target keep is
  the stored path of the source keep. So stored paths are natural in codes.
\item \ag{cohA}: for keeps \ag{δ} (with code \ag{q}) and \ag{σ}
  (with code \ag{e}), the two stored paths compose to the stored path of the
  composite keep (with code \ag{e'}), up to the composition path.
\item \ag{cohW}: here $U = \mathrm{unwk}$, $\delta^+ =
  \ag{wkThin δ p}$, $\bar a = \mathrm{unwk}\,a$, and $u(\delta)$ is
  \ag{unwk-nat δ}. The image under $U$ of the stored path of a weakened keep
  is the stored path at $\bar a$, up to the weakening comparisons.
\end{itemize}
This is where the naive interpretation used UIP: it needed two transports to agree,
and now we need two paths to be equal, but only the specific ones that arise.

\subsection{The 2-cells}

The three stored-path cells are not proved directly. They are pasted from a
small set of \emph{generic} 2-cells, which are proved by induction on the
syntax. There are six generic 2-cells. We write $\circ(\delta,\sigma)$ for \ag{⟦∘⟧ δ σ a}, and $u(\delta)$,
$U$, $\delta^+$, $\bar a$ and $\langle c\rangle$ as before; each cell says
that its diagram commutes.
\begin{center}
\begin{tikzcd}[column sep=large,row sep=large]
\sigma a \ar[r,"\langle c\rangle"] \ar[dr,"\langle cd\rangle"'] & \tau a \ar[d,"\langle d\rangle"] \\
\phantom{x} & \rho a
\ar[from=1-1,to=2-2,phantom,"\text{\ag{cohT}}" pos=0.55,shift left=2.5ex]
\end{tikzcd}
\qquad
\begin{tikzcd}[column sep=large,row sep=large]
(\delta\circ\sigma)a \ar[r,"{\circ(\delta,\sigma)}"] \ar[d,"\langle c'\rangle"'] & \delta\sigma a \ar[d,"\delta\langle c\rangle"] \\
(\delta\circ\tau)a \ar[r,"{\circ(\delta,\tau)}"'] & \delta\tau a
\ar[from=1-1,to=2-2,phantom,"\text{\ag{cohR}}"]
\end{tikzcd}
\qquad
\begin{tikzcd}[column sep=large,row sep=large]
(\delta\circ\sigma)a \ar[r,"{\circ(\delta,\sigma)}"] \ar[d,"\langle c'\rangle"'] & \delta\sigma a \ar[d,"\langle c\rangle"] \\
(\delta_2\circ\sigma)a \ar[r,"{\circ(\delta_2,\sigma)}"'] & \delta_2\sigma a
\ar[from=1-1,to=2-2,phantom,"\text{\ag{cohL}}"]
\end{tikzcd}

\medskip
\begin{tikzcd}[column sep=large,row sep=large]
(\rho\circ(\delta\circ\sigma))a \ar[r,"{\circ(\rho,\delta\circ\sigma)}"] \ar[d,"\langle c\rangle"'] & \rho((\delta\circ\sigma)a) \ar[r,"{\rho\circ(\delta,\sigma)}"] & \rho\delta\sigma a \\
((\rho\circ\delta)\circ\sigma)a \ar[r,"{\circ(\rho\circ\delta,\sigma)}"'] & (\rho\circ\delta)\sigma a \ar[ur,"{\circ(\rho,\delta)}"'] &
\ar[from=1-1,to=2-2,phantom,"\text{\ag{cohAs}}"]
\end{tikzcd}
\qquad
\begin{tikzcd}[column sep=large,row sep=large]
U(\rho^+a) \ar[r,"U\langle c'\rangle"] \ar[d,"u(\rho)"'] & U(\rho'^+a) \ar[d,"u(\rho')"] \\
\rho\bar a \ar[r,"\langle c\rangle"'] & \rho'\bar a
\ar[from=1-1,to=2-2,phantom,"\text{\ag{cohWN}}"]
\end{tikzcd}

\medskip
\begin{tikzcd}[column sep=large,row sep=large]
U((\delta^+\circ\sigma^+)a) \ar[r,"U\langle c\rangle"] \ar[dr,"{U\circ(\delta^+,\sigma^+)}"'] & U((\delta\circ\sigma)^+a) \ar[r,"{u(\delta\circ\sigma)}"] & (\delta\circ\sigma)\bar a \ar[r,"{\circ(\delta,\sigma)}"] & \delta\sigma\bar a \\
\phantom{x} & U(\delta^+\sigma^+a) \ar[r,"{u(\delta)}"'] & \delta\,U(\sigma^+a) \ar[ur,"{\delta\,u(\sigma)}"'] &
\ar[from=1-2,to=2-3,phantom,"\text{\ag{cohW∘}}"]
\end{tikzcd}
\end{center}
Here $c, c', cd, \dots$ are codes between the relevant thinnings. In words: code paths compose (\ag{cohT}),
composition is natural in codes on the right and on the left (\ag{cohR},
\ag{cohL}), composition is associative up to a code (\ag{cohAs}), and the
weakening comparison is natural in codes and compatible with composition
(\ag{cohWN}, \ag{cohW∘}). So the semantics has the structure of a
pseudofunctor on the category of thinnings: \ag{≈} is its action on codes,
\ag{⟦∘⟧} its compositor, \ag{unwk-nat} a pseudonatural comparison, and the
generic 2-cells are the corresponding naturality and associativity data.

\paragraph{The stored-path cells as pastings.}
Each stored-path cell is a pasting of generic 2-cells
(Figure~\ref{fig:cells}). The pasting is done by a lemma of pure path
algebra, proved by J, since it holds for arbitrary paths. For \ag{cohN} it is
\begin{agda}
lemN : m ∙ n ≡ k ∙ l → k ∙ e₂ ≡ e₁ → sym n ∙ (sym m ∙ e₁) ≡ sym l ∙ e₂
\end{agda}
instantiated with $m = \circ(\delta,\sigma)$, $n = \delta\langle c\rangle$,
$k = \langle c'\rangle$, $l = \circ(\delta,\tau)$, $e_1 = \langle e\rangle$
and $e_2 = \langle e'\rangle$. Then the first hypothesis is the \ag{cohR}
square for $c$ and the induced code $c'$, the second is the \ag{cohT}
triangle for $c'$, $e'$ and $e$, and since the stored path is
(writing, as in Agda, $p\pc q$ for path composition)
$\mathrm{kp}_e = \circ(\delta,\sigma)^{-1}\pc\langle e\rangle$, the conclusion
$\delta\langle c\rangle^{-1}\pc\mathrm{kp}_e = \mathrm{kp}_{e'}$ is
\ag{cohN}: a square and a triangle pasted along $\langle c'\rangle$
(Figure~\ref{fig:cells}(a)). In the same way \ag{lemA} pastes \ag{cohAs},
\ag{cohL} and two instances of \ag{cohT} into \ag{cohA}, and \ag{lemW} pastes
\ag{cohW∘}, \ag{cohWN} and \ag{cohT} into \ag{cohW}.

\begin{figure}[!htbp]
\centering
\small
\textbf{(a)} \ag{cohN}\\[1ex]
\begin{tikzcd}[column sep=large,row sep=large]
\delta\sigma a \ar[d,"\delta\langle c\rangle"'] & (\delta\circ\sigma)a \ar[l,"{\circ(\delta,\sigma)}"'] \ar[d,"\langle c'\rangle"] \ar[dr,"\langle e\rangle"] & \phantom{x} \\
\delta\tau a & (\delta\circ\tau)a \ar[l,"{\circ(\delta,\tau)}"] \ar[r,"\langle e'\rangle"'] & \delta' a
\ar[from=1-1,to=2-2,phantom,"\mathsf R"]
\ar[from=2-2,to=1-3,phantom,"\mathsf T" pos=0.3]
\end{tikzcd}

\bigskip
\textbf{(b)} \ag{cohA}\\[1ex]
\begin{tikzcd}[column sep=large,row sep=large]
\phantom{x} & (\delta_C\circ(\delta\circ\sigma))a \ar[r,"{\circ(\delta_C,\delta\circ\sigma)}"] \ar[d,"\langle c_1\rangle"] \ar[dl,"\langle e'\rangle"']
  & \delta_C((\delta\circ\sigma)a) \ar[r,"{\delta_C\circ(\delta,\sigma)}"] & \delta_C\delta\sigma a \\
\delta_A a & ((\delta_C\circ\delta)\circ\sigma)a \ar[r,"{\circ(\delta_C\circ\delta,\sigma)}"] \ar[d,"\langle c_2\rangle"] \ar[l,"\langle c_2e\rangle"]
  & (\delta_C\circ\delta)\sigma a \ar[ur,"{\circ(\delta_C,\delta)}"'] \ar[d,"\langle q\rangle"] & \\
\phantom{x} & (\delta_B\circ\sigma)a \ar[r,"{\circ(\delta_B,\sigma)}"'] \ar[ul,"\langle e\rangle"] & \delta_B\sigma a &
\ar[from=2-2,to=1-3,phantom,"\mathsf{As}"]
\ar[from=2-2,to=3-3,phantom,"\mathsf L"]
\ar[from=2-2,to=1-1,phantom,"\mathsf T" pos=0.35]
\ar[from=2-2,to=3-1,phantom,"\mathsf T" pos=0.35]
\end{tikzcd}

\bigskip
\textbf{(c)} \ag{cohW}\\[1ex]
\begin{tikzcd}[column sep=large,row sep=large]
U((\delta^+\circ\sigma^+)a) \ar[r,"{U\circ(\delta^+,\sigma^+)}"] \ar[d,"U\langle c_3\rangle"] \ar[dd,bend right=75,looseness=1.6,"U\langle e_1\rangle"']
  & U(\delta^+\sigma^+a) \ar[r,"{u(\delta)}"] & \delta\,U(\sigma^+a) \ar[r,"{\delta\,u(\sigma)}"] & \delta\sigma\bar a \\
U((\delta\circ\sigma)^+a) \ar[rrr,"{u(\delta\circ\sigma)}"] \ar[d,"U\langle c_w\rangle"] & & & (\delta\circ\sigma)\bar a \ar[u,"{\circ(\delta,\sigma)}"'] \ar[d,"\langle e\rangle"] \\
U(\delta'^+a) \ar[rrr,"{u(\delta')}"'] & & & \delta'\bar a
\ar[from=1-1,to=2-4,phantom,"\mathsf{W\circ}"]
\ar[from=2-1,to=3-4,phantom,"\mathsf{WN}"]
\ar[from=1-1,to=3-1,phantom,"\mathsf T",shift right=17mm]
\end{tikzcd}

\caption{The stored-path cells as pastings of the generic 2-cells. We omit
$\llbracket-\rrbracket$: $\delta\sigma a$ stands for
$\llbracket\delta\rrbracket(\llbracket\sigma\rrbracket\,a)$ and
$(\delta\circ\sigma)a$ for $\llbracket\delta\circ\sigma\rrbracket\,a$, and,
as in Agda, $p\pc q$ for path composition. A
stored path is $\mathrm{kp}\,\delta\,\sigma\,\delta'\,e =
\circ(\delta,\sigma)^{-1}\pc \langle e\rangle$. In each diagram the cell is
the equality of the two paths around the outer boundary, from
$\delta\sigma a$ in (a) and from the top right corner in (b) and (c), and it is obtained by pasting the labelled regions; this is what the
J-lemmas \ag{lemN}, \ag{lemA}, \ag{lemW} do.
(a) $\mathrm{kp}\,\delta\,\sigma\,\delta'\,e = \delta\langle c\rangle\pc 
\mathrm{kp}\,\delta\,\tau\,\delta'\,e'$, where $c'$ is the code induced by $c$.
(b) $\delta_C\circ(\delta,\sigma)^{-1}\pc \mathrm{kp}\,\delta_C\,(\delta\circ\sigma)\,\delta_A\,e'
= \mathrm{kp}\,\delta_C\,\delta\,\delta_B\,q\pc \mathrm{kp}\,\delta_B\,\sigma\,\delta_A\,e$,
where $c_1$ is the associativity code, $c_2$ the code induced by $q$, and $c_2e$
their composite with $e$.
(c) $U = \mathrm{unwk}$, $\delta^+ = \mathrm{wkThin}\,\delta\,p$ and $\bar a =
\mathrm{unwk}\,a$: the image under $U$ of the stored path of the weakened keep
agrees with the stored path at $\bar a$, up to the weakening comparisons $u$.}
\label{fig:cells}
\end{figure}

\paragraph{Proving the generic 2-cells.}
The generic 2-cells are proved by induction on the thinnings, mutually with
the interpretation. On \ag{tε} and \ag{t-d} they reduce to the induction
hypothesis or hold by \ag{refl}. On keeps both sides are composites of
\ag{pathK}s, and two such composites are equal if their base paths are equal
(the induction hypothesis) and their stored components agree. The stored
components are 2-paths between paths in \ag{⟦ C ⟧T γ}, so their agreement is
a \emph{3-path} in \ag{⟦ C ⟧T γ}.

\subsection{Where the groupoids are used}\label{sec:where}

This is the only place where we use that the semantic types are groupoids. In a
groupoid the path types are sets, so any two parallel 2-paths are equal and the
required 3-paths exist:
\begin{agda}
grpdEq : isGroupoid D → {u v : D}{q q' : u ≡ v}(s s' : q ≡ q') → s ≡ s'
\end{agda}
This is used in the two lemmas that compare composites of \ag{pathK}s
(\ag{pathK-∙∙} and \ag{pathK-∙≡}), and nowhere else. So the tower of
coherences stops after one step:
\begin{center}
\begin{tabular}{lll}
\toprule
level & cells & proved by\\
\midrule
1 & \ag{≈}, \ag{⟦∘⟧}, \ag{unwk-nat} & recursion; stored parts \ag{cohN}, \ag{cohA}, \ag{cohW}\\
2 & \ag{cohT}, \ag{cohR}, \ag{cohL}, \ag{cohAs}, \ag{cohWN}, \ag{cohW∘} & induction; stored parts 3-paths\\
3 & 3-paths in the semantic types & trivial in groupoids\\
\bottomrule
\end{tabular}
\end{center}
The naive interpretation needed the analogous statement one level
lower, that 2-paths in the semantic types are trivial, which is UIP. The
Forded semantics and the action of codes shift the use of truncation up by
one dimension.

\paragraph{The formalisation.}
\agfile{experiments/sem1/Sem1.agda} has about 870 lines, contains no
postulates and no holes, and uses no K (it is checked with
\ag{--cubical}, which implies \ag{--without-K}). It uses 15
\ag{TERMINATING} pragmas: the mutual recursion is structural, but it goes
through the decoded codes and through the mutually defined functions in a way
that Agda's termination checker does not recognise. Type checking takes about
two and a half minutes. Most of the work is in the keep cases, where the J-lemmas
must be applied with explicit motives, and in the cases of \ag{cohAs} and
\ag{cohW∘} that combine several keeps.

\section{Semi-simplicial groupoids}\label{sec:sstgrpd}

Interpreting the uniform construction of Section~\ref{sec:tt} gives
\begin{agda}
⟦ SST n ⟧C : Set₁
\end{agda}
for a variable $n$: a type of semi-simplicial groupoids up to level $n$, in
Forded form. Each level is a family of groupoids over the Forded telescope
of the boundary of a simplex, whose entries store paths from the selected
faces to their points.

The file \agfile{experiments/sem1/SemSST.agda} checks that the
construction selects the right faces. For $n = 1, 2$ there are maps from
the standard, unforded semi-simplicial groupoids, written out by hand, into
$\llbracket\mathrm{SST}\,n\rrbracket$, which transport each value along its
stored path to the family applied to the faces it should select. These maps
typecheck only if every face map of the construction selects the right lower
simplices. For $n = 3$ the same check is too expensive, since Agda would have to
normalise stored paths that contain the coherence proofs. But since a
thinning never changes the point $d$ or the value $v$ of an entry, it is
enough to check that each triangle face of a tetrahedron selects entries with
the right points and values. These checks hold by \ag{refl}, and a
deliberately wrong face fails, as expected.

We have not formalised that $\llbracket\mathrm{SST}\,n\rrbracket$ is
equivalent to the type of semi-simplicial groupoids up to level $n$ written
out by hand, for each numeral $n$. Level by level this is the contraction of
the singletons $(d, q)$, but the equivalence has to be compatible with the
face maps, and we have only checked this for $n \le 3$ as above.

The natural statement of correctness uses two-level type
theory~\cite{ack-strict,acks-2ltt}, where semi-simplicial types up to level
$n$ can be defined for an outer natural number $n$. Restricting that
definition to families of groupoids gives a type
$\mathrm{SST}^{\mathrm{Grpd}}_n$ of semi-simplicial groupoids, and we expect
\[
\llbracket\mathrm{SST}\,n\rrbracket \simeq \mathrm{SST}^{\mathrm{Grpd}}_n
\]
for every outer $n$. Since our construction of $\mathrm{SST}\,n$ follows the
recursion of~\cite[Section~3]{ack-strict}, the proof would go by induction
along it. The skeleta are equivalent by contracting the singletons $(d,q)$,
and the main work is to show that these equivalences commute with the face
maps, coherently with $\alpha$: the face maps compose strictly in two-level
type theory, and only up to the 2-cells of Section~\ref{sec:grpd} in our
semantics. We leave this for future work.

\section{The next level}\label{sec:higher}

The same construction, with groupoids replaced by 2-types, is in
\agfile{experiments/sem1/Sem2.agda}. Now the stored components of the generic
2-cells are 3-paths, which are no longer trivial. Every place where
\ag{Sem1} used \ag{grpdEq} becomes an obligation: a 3-path between two
specific pastings of 2-cells, transported along a 2-cell at smaller
thinnings. There are six such obligations, one for each generic 2-cell, and
unfolding the stored-path cells turns each into an equation between two
pastings of the generic 2-cells, that is, a 3-cell. We expect twelve generic
3-cells:
\begin{center}
\begin{tabular}{ll}
\toprule
3-cell & statement\\
\midrule
\ag{T3} & \ag{cohT} is associative\\
\ag{RT}, \ag{LT} & \ag{cohR}, \ag{cohL} are functorial in codes\\
\ag{RL} & interchange of \ag{cohR} and \ag{cohL}\\
\ag{AsR}, \ag{AsM}, \ag{AsL} & \ag{cohAs} is natural in each argument\\
\ag{Pent} & the pentagon for \ag{cohAs}\\
\ag{WNT} & \ag{cohWN} is functorial in codes\\
\ag{W∘R}, \ag{W∘L} & \ag{cohW∘} is natural in each argument\\
\ag{W∘As} & \ag{cohW∘} is compatible with \ag{cohAs}\\
\bottomrule
\end{tabular}
\end{center}
These are the pseudofunctor axioms one level up. So far we have discharged the
obligation for \ag{cohT} from \ag{T3} and \ag{RT}, and proved \ag{T3}. On keeps
the proof of \ag{T3} again splits into a base 3-cell, the induction hypothesis,
and a stored component, which is a 4-path in a 2-type and therefore trivial.
This confirms the pattern: at each level the cells are equations between two
given parallel pastings of the cells below, so no boundary data is ever
needed, and the keep cases are discharged by the induction hypothesis together
with the truncation one dimension up. The remaining obligations and 3-cells
are work in progress.

\paragraph{The pattern.}
The cells are indexed by simplices of the nerve of the category of
thinnings. The compositor $\circ(\delta,\sigma)$ belongs to a 2-simplex with
edges $\sigma$, $\delta$, $\delta\circ\sigma$; the associator \ag{cohAs} to a
3-simplex, whose four faces are compositor cells; the pentagon to a
4-simplex, with five associator faces. The code cells express that the
semantics respects the equations of the syntax, and the weakening cells that
it commutes with context extension. To define the interpretation into
$k$-types uniformly in an internal $k:\mathbb N$, one would have to describe
this whole tower as a single statement, ``parallel pastings of the
generators agree'', generated by induction on the dimension. Stating it needs
a uniform notion of pasting, and this is where higher categories come back:
the coherences say that the interpretation is a $k$-coherent functor from the
thinnings into the universe of $k$-types, and the latter is described by
semi-simplicial types of higher dimension.

\section{Discussion}\label{sec:discussion}

\paragraph{What has been achieved.}
The obstruction of~\cite{obstruction} is not a limitation of the
syntax, nor of the construction of $\mathrm{SST}\,n$ in it. It is the
coherence of the semantic composition, which the naive interpretation takes
from UIP. With a Forded semantics of extension and a semantic action of
codes, these coherences can be proved by induction, and for groupoids the
induction closes after one level. For 2-types it needs one more level, and the
shape of that level is now known.

\paragraph{Induction-recursion.}
The construction relies on the inductive-recursive syntax: composition and
weakening are defined by recursion and occur in the indices of constructors,
so that the syntax has no transports and its recursive equations hold
definitionally. There are general accounts of
induction-recursion~\cite{dybjer-ir,dybjer-setzer} and of
induction-induction~\cite{forsberg-setzer}, but we know of none that covers this
form. In~\cite{obstruction} we also considered syntaxes without
induction-recursion: a Forded higher inductive-inductive type, which is a set
but whose naive model needs UIP, and which could instead be interpreted
through normal forms, that is, through the inductive-recursive syntax. A
different possibility is a Segal-style syntax in which composition is a horn
filler with a contraction, and the thinnings are truncated at the level of the
semantics (file \agfile{experiments/segal/TT1.agda}, a higher
inductive-inductive type at the level of groupoids). An interpretation of that
syntax in groupoids would avoid induction-recursion altogether, and would
organise the coherences as the Segal condition for the universe. We have not
yet attempted it.

\paragraph{Higher categories.}
The pattern of Section~\ref{sec:higher} suggests organising the
construction with higher categories. Instead of extracting the coherences of
each level from the syntax, we would define, mutually and by induction on the
truncation level $k$, semi-simplicial $k$-types and $(k,1)$-categories
presented by them, and use the higher-categorical structure of the universe
of $k$-types to interpret the syntax at the next level. The coherences of
Sections~\ref{sec:grpd} and~\ref{sec:higher} would then be instances of this
structure rather than ad hoc lemmas. We think this is feasible, and
conceptually the right picture: truncated semi-simplicial types and higher
categories should be definable together. Passing from the truncated
versions to untruncated semi-simplicial types is a separate step. If the
types $\mathrm{SST}^{(k)}_n$ of semi-simplicial $k$-types are definable
uniformly in $k$, together with maps
$\mathrm{SST}^{(k+1)}_n \to \mathrm{SST}^{(k)}_n$, we can simply
\emph{define}
\[
\mathrm{SST}_n = \lim_k \mathrm{SST}^{(k)}_n,
\]
the type of Postnikov towers of truncated semi-simplicial types. The maps
are not levelwise truncation: a truncated family has to live over the
boundary built from the truncated lower levels, so they are the Postnikov
truncation of the diagram (truncate the total spaces of simplices, then take
fibres over the new boundaries), and defining them uniformly in $n$ needs the
semantics to be natural in the truncation level. Whether this limit is
\emph{the} type of semi-simplicial types depends on whether types are
determined by their truncations (Postnikov towers converging). This is where
hypercompleteness comes in (Whitehead's principle: $\infty$-connected maps
are equivalences), which is not provable in book HoTT. So the outcome would be
that semi-simplicial types are definable up to Postnikov completion, and
agree with the intended ones in models where Postnikov towers converge, such
as simplicial sets. It would also suggest why a direct definition of the
untruncated type is hard to find: it would have to work in models where
towers do not converge.

\section*{Use of AI tools}
This note and the accompanying Agda development were produced in
collaboration with Claude (Anthropic), used through Claude Code. Claude
drafted text and Agda code, carried out experiments in Agda, and searched the
literature, under the author's direction. The author has checked the content
and is responsible for it.

\begin{thebibliography}{99}
\bibitem{ford}
Henry Ford, in collaboration with Samuel Crowther.
\newblock \emph{My Life and Work}.
\newblock Doubleday, Page \& Company, 1922.

\bibitem{obstruction}
Thorsten Altenkirch.
\newblock Any colour you like, as long as it's UIP: semi-simplicial types from
  weakenings and thinnings.
\newblock arXiv:2610.29585 [cs.LO], 2026.
  \href{https://doi.org/10.48550/arXiv.2610.29585}{doi:10.48550/arXiv.2610.29585}.

\bibitem{ack-strict}
Thorsten Altenkirch, Paolo Capriotti and Nicolai Kraus.
\newblock Extending homotopy type theory with strict equality.
\newblock In \emph{CSL 2016}, LIPIcs 62, pages 21:1--21:17, 2016.

\bibitem{acks-2ltt}
Danil Annenkov, Paolo Capriotti, Nicolai Kraus and Christian Sattler.
\newblock Two-level type theory and applications.
\newblock \emph{Mathematical Structures in Computer Science}, 33(8):688--743,
  2023.

\bibitem{dybjer-ir}
Peter Dybjer.
\newblock A general formulation of simultaneous inductive-recursive
  definitions in type theory.
\newblock \emph{Journal of Symbolic Logic}, 65(2):525--549, 2000.

\bibitem{dybjer-setzer}
Peter Dybjer and Anton Setzer.
\newblock A finite axiomatization of inductive-recursive definitions.
\newblock In \emph{TLCA 1999}, LNCS 1581, pages 129--146. Springer, 1999.

\bibitem{hott-book}
The Univalent Foundations Program.
\newblock \emph{Homotopy Type Theory: Univalent Foundations of Mathematics}.
\newblock Institute for Advanced Study, 2013.

\bibitem{dtt}
Astra Kolomatskaia and Michael Shulman.
\newblock Displayed type theory and semi-simplicial types.
\newblock \emph{Mathematical Structures in Computer Science}, published online
  2025. arXiv:2311.18781.

\bibitem{forsberg-setzer}
Fredrik Nordvall Forsberg and Anton Setzer.
\newblock Inductive-inductive definitions.
\newblock In \emph{CSL 2010}, LNCS 6247, pages 454--468. Springer, 2010.

\bibitem{shulman-eat}
Michael Shulman.
\newblock Homotopy type theory should eat itself (but so far, it's too big to
  swallow).
\newblock Blog post, \url{https://homotopytypetheory.org/2014/03/03/hott-should-eat-itself/},
  2014.

\bibitem{ufias-sst}
Univalent Foundations special year, IAS.
\newblock Semi-simplicial types (open problem page).
\newblock \url{https://ncatlab.org/ufias2012/show/Semi-simplicial+types},
  2012--13.

\bibitem{cubical-agda}
Andrea Vezzosi, Anders Mörtberg and Andreas Abel.
\newblock Cubical Agda: a dependently typed programming language with
  univalence and higher inductive types.
\newblock \emph{Journal of Functional Programming}, 31:e8, 2021.
\end{thebibliography}

\end{document}
